% =====================================================================
%  30.2  Cayley's formula and the forest-completion variant.
%  Exports: sec:c32-cayley  thm:c32-cayley  thm:c32-rooted
%  thm:c32-forest  lem:c32-forestden  def:c32-prufv
%  prob:c32-prufvtx  fig:c32-prufvtx  lst:c32-prufvtx
%  lst:c32-modpow
%  Solved CF 2190D "Prufer Vertex" (2500, L5), live-verified against the CF
%  problem page (title "D. Prufer Vertex", 2s TL, 256 MB, t<=1e4, n<=2e5,
%  m<=n-1, sum n<=2e5, mod 998244353).  Three official samples:
%  case 1 (n=3, m=0, forest of three isolates) -> "1 2"; case 2 (n=5, m=4,
%  tree already) -> "0 0 0 1"; case 3 (n=6, m=3, forest of 3 components of
%  sizes 4+1+1) -> "12 0 1 6 5", matching n^{k-2} * prod s_i = 6 * 4 = 24.
% =====================================================================
\section{Cayley's Formula and the Forest-Completion Count}\label{sec:c32-cayley}

\begin{previouslybox}[Ch.30.1 turned trees into sequences]
\(\secref{c32-prufer}\) established the bijection \(\mathcal{T}_n
\leftrightarrow [n]^{n-2}\).  Sequences are easy to count.  This section
cashier's that bijection against the two enumerative workhorses every later
section uses: Cayley's formula \(n^{n-2}\), and its \emph{forest-completion}
variant \(n^{k-2}\prod_i s_i\), which counts the ways to add \(k-1\) edges to
a \(k\)-component forest to obtain one tree.
\end{previouslybox}

\begin{topicintro}
The Pr\"ufer code is a sequence of length \(n-2\) with entries in \([n]\); the
count of such sequences is \(n^{n-2}\), and each corresponds to exactly one
labelled tree (\(\thmref{c32-bij}\) from \(\secref{c32-prufer}\)).  That is
Cayley's formula (\(\thmref{c32-cayley}\)) in one line.  Rooting multiplies by
\(n\) (\(\thmref{c32-rooted}\)); prescribing degrees gives the multinomial
\(\corref{c32-mult}\).  The interesting variant is: given a forest \(F\) on
\([n]\) with \(k\) components of sizes \(s_1,\dots,s_k\), count the number of
labelled trees \(T\in\mathcal{T}_n\) that \emph{contain} \(F\) (add edges,
never remove).  The answer is \(n^{k-2}\prod_i s_i\) (\(\thmref{c32-forest}\)),
which we prove by a Pr\"ufer-style count on contracted components rather than
by an independent bijection -- the same tool does the work.

CF 2190D (\(\probref{c32-prufvtx}\)) is the summit of \(\secref{c32-cayley}\):
read \(P(T)\), the "second survivor" of the smallest-leaf peel on \(T\), and
ask for a histogram of \(P(T)\) across all completions of a given forest.  It
uses \(\thmref{c32-forest}\) and a per-vertex code-structure observation to
turn the count into a fold that runs in \(\bigO{n\log n}\) per test.

\textbf{Prerequisites.} \(\secref{c32-prufer}\) throughout -- the bijection,
the degree-from-multiplicity lemma, the multinomial.  Ch.~17 (\chref{c28}) --
DSU or the DFS equivalent for reading the forest's \(k\) and \(s_i\).
Modular inverse and fast exponentiation (Ch.~6 background), used only through
standard \(\mathrm{pow}(a,b)\bmod p\).

\textbf{After this section you can}
\begin{itemize}[nosep]
  \item state Cayley \(n^{n-2}\) as a corollary of \(\thmref{c32-bij}\); \hfill\lv{3}
  \item count completions of a forest via \(n^{k-2}\prod s_i\); \hfill\lv{4}
  \item read a labelled tree's smallest-leaf survivors and identify \(P(T)\); \hfill\lv{4}
  \item solve CF 2190D: for each \(v<n\), count completions whose \(P(T)=v\). \hfill\lv{5}
\end{itemize}
\end{topicintro}

% ---------------------------------------------------------------
\subsection{Cayley's formula, rooted trees, and the degree multinomial}
% ---------------------------------------------------------------

\begin{theorem}{Cayley}{c32-cayley}
For every \(n\ge 1\), \(|\mathcal{T}_n| = n^{n-2}\).
\end{theorem}

\begin{proofbox}
\(\thmref{c32-bij}\) puts \(\mathcal{T}_n\) in bijection with the set of
length-\((n-2)\) sequences over \([n]\), which has size \(n^{n-2}\) by the
rule of product.  Boundary \(n=1\): \(1^{-1}\) reads as \(1\), the unique
one-vertex tree.  Boundary \(n=2\): \(2^0=1\), the single edge \(1\!-\!2\).
\(\Box\)
\end{proofbox}

\begin{theorem}{Rooted labelled trees}{c32-rooted}
The number of pairs \((T,r)\) with \(T\in\mathcal{T}_n\) and \(r\in[n]\) the
chosen root is \(n^{n-1}\).
\end{theorem}

\begin{proofbox}
Multiply Cayley by \(n\): each tree admits \(n\) distinct rootings. \(\Box\)
\end{proofbox}

\begin{lstlisting}[caption={Modular exponentiation, used to raise \(n\) to a fixed exponent modulo \(p\) in every counting application of this section. Iterative binary power, \(\bigO{\log e}\)}, label={lst:c32-modpow}]
#include <bits/stdc++.h>
using namespace std;
typedef long long ll;
ll modPow(ll base, ll exp, ll mod){
  ll res=1; base%=mod;
  while(exp>0){
    if(exp&1) res=(res*base)%mod;
    base=(base*base)%mod;
    exp>>=1;
  }
  return res;
}
\end{lstlisting}

% ---------------------------------------------------------------
\subsection{The forest-completion formula}
% ---------------------------------------------------------------

\begin{definition}{Completion of a forest}{c32-completion}
Let \(F=(V,E_F)\) be a forest on \(V=[n]\) with \(k\) connected components of
sizes \(s_1,\dots,s_k\) (so \(\sum_i s_i = n\)).  A \term{completion} of \(F\)
is a labelled tree \(T\in\mathcal{T}_n\) with \(E_F\subseteq E(T)\); adding
\(k-1\) edges to \(F\) turns it into a tree.
\end{definition}

\begin{theorem}{Forest-completion count}{c32-forest}
The number of completions of \(F\) is \(n^{k-2}\prod_{i=1}^{k} s_i\).  In
particular, when \(k=n\) (empty forest, all \(s_i=1\)) this recovers Cayley;
when \(k=1\) (already a tree) it gives \(1\).
\end{theorem}

\begin{proofbox}
Two proofs exist.  The classical one contracts each component to a
"super-vertex" and applies the Pr\"ufer code to the resulting tree of
components, weighting each super-vertex by the size \(s_i\) of its
component -- the code has length \(k-2\), each entry is a super-vertex, and
each position of the code is expanded in \(s_i\) ways to a real label.  A
more direct proof uses the Matrix-Tree theorem (\(\secref{c32-matrix}\)) on
the graph \(K_n\) with \(E_F\) forced to weight \(0\) and every other edge to
weight \(1\); the Laplacian cofactor is a block determinant that evaluates to
\(n^{k-2}\prod s_i\).  We sketch the Pr\"ufer version here since it is the
one that generalises to \(\probref{c32-prufvtx}\).

\emph{Pr\"ufer-style proof.}  Given a completion \(T\), prune \(T\) down by
repeatedly removing \emph{real} leaves belonging to nontrivial components;
when a component's vertices are all gone except its "root," the component
itself becomes a leaf of the residual component-tree.  The resulting
component-code has length \(k-2\) and entries in \(\{1,\dots,k\}\), i.e.\
\(k^{k-2}\) possibilities.  Now, in the original tree \(T\), the position
\(j\) of this component-code records which component is entered from outside;
pick the specific label within the entered component in \(s_{\text{entered}}\)
ways, and pick the specific label within the exiting component in
\(s_{\text{exiting}}\) ways.  Multiplying these local choices over the
\(k-1\) component-edges and combining with the \(k^{k-2}\)-way Pr\"ufer
arrangement yields \(n^{k-2}\prod s_i\) after a short algebraic identity
(the "double count" that appears in most textbook proofs of this theorem).
\(\Box\)
\end{proofbox}

\begin{insightbox}[Reading the formula]
The \(n^{k-2}\) factor is the Cayley count on a "contracted graph" of \(k\)
components weighted to sum \(n\); the \(\prod s_i\) factor is the local
choice of \emph{which} vertex in each component is used to graft the new
edge.  Both factors grow polynomially in \(n\), so a single fast-exponentiation
call plus one loop over the DSU census is the whole computation.
\end{insightbox}

% ---------------------------------------------------------------
\subsection{Solved: CF 2190D -- Prufer Vertex (2500, L5)}
% ---------------------------------------------------------------

\solvedproblem{c32-prufvtx}{Prufer Vertex}{CF 2190D}{https://codeforces.com/problemset/problem/2190/D}{CF 2500 L5}{Pruefer code (\(\secref{c32-prufer}\)), forest-completion count (\(\thmref{c32-forest}\)), tree-DP fold along the component spine, modular arithmetic}

\begin{problemstatement}[CF 2190D, verified against the live CF problem page]
You are given a forest on \([n]\) with \(m\) edges; the forest has \(k\)
connected components of sizes \(s_1,\dots,s_k\).  There are exactly
\(n^{k-2}\prod_i s_i\) ways to add edges to the forest so that it becomes a
single labelled tree on \([n]\).  For a labelled tree \(T\) with \(n\ge 2\),
define its \term{Prufer vertex} \(P(T)\) as follows: peel the tree using the
standard smallest-leaf prune until only two vertices remain (one of them is
always the label \(n\)); \(P(T)\) is the \emph{other} remaining vertex.  For
each \(v\) with \(1\le v<n\), count the completions \(T\) of the given forest
with \(P(T)=v\), modulo \(998244353\).  Constraints: \(t\le 10^4\) tests,
\(n\le 2\cdot 10^5\), \(m\le n-1\), sum of \(n\) over all tests \(\le
2\cdot 10^5\).  Samples (three cases): \(n=3, m=0\) outputs "1 2"; \(n=5,
m=4\) (already a tree) outputs "0 0 0 1"; \(n=6, m=3\) outputs "12 0 1 6 5".
\end{problemstatement}

\begin{figure}[H]
\centering
\begin{tikzpicture}[diagram, scale=0.9, transform shape,
   tnode/.style={circle, draw=primaryblue!70, fill=primaryblue!8, line width=0.7pt,
                 minimum size=7mm, font=\scriptsize\bfseries, inner sep=1pt},
   comp/.style={rectangle, rounded corners=2pt, draw=difforange!70, fill=difforange!10,
                line width=0.8pt, minimum width=14mm, minimum height=9mm,
                font=\scriptsize\bfseries, inner sep=2pt}]
  \node[comp] (cA) at (0,0) {\(\{1,2,4,6\}\)};
  \node[comp] (cB) at (3.2,0) {\(\{3\}\)};
  \node[comp] (cC) at (5.2,0) {\(\{5\}\)};
  \draw[difforange!70, line width=1pt] (cA)--(cB)--(cC);
  \node[font=\scriptsize, text width=6cm, align=center] at (2.4,-1.3)
    {case 3: three components, sizes \(4,1,1\); completions \(=6^{3-2}\cdot 4\cdot 1\cdot 1=24\)};
\end{tikzpicture}
\caption{Case 3 of the sample.  The forest has three components; \(24\)
completions distribute as \(P(T)\)-values \(12,0,1,6,5\) across
\(v=1,\dots,5\), summing back to \(24\) as required.}
\label{fig:c32-prufvtx}
\end{figure}

\begin{attemptgate}[Attempt it yourself first -- \lv{5}, give yourself 45 minutes]
Before coding: what structural property of the smallest-leaf peel tells you
\(P(T)\) directly, without simulating the whole peel on each completion?
\peek{The peel is exactly Pr\"ufer encoding, so \(P(T)\) is one of the two
labels still in the decode heap when the code is exhausted -- specifically the
one that is \emph{not} \(n\).  Characterising "\(v\) is in the final heap" is a
constraint on the code's residual multiset.}
\nudge{For each completion \(T\), \(P(T)=v\) iff, in the Pr\"ufer code of \(T\),
the label \(v\) is the largest label that either appears in the code OR, when
it does not, is not beaten as the smallest leaf by \(n\) until the last step.
Read \(\lemref{c32-deg}\) and \(\thmref{c32-forest}\) together: a completion
\(T\) of \(F\) is a tree on \([n]\) containing \(E_F\), and its code length is
\(n-2\) with degrees constrained by the DSU component sizes.}
\fullhint{Root each component of \(F\) at its largest label.  Then \(P(T)\) for
a completion \(T\) can be characterised by a fold on the component-tree
formed by contracting each component to its root: the fold computes, per
component root, the number of completions whose second-to-last surviving
vertex lies inside the component.  Combine across components using
\(\thmref{c32-forest}\) as a base case; a per-vertex count is then a
convolution of subtree polynomials, in \(\bigO{n\log n}\) per test with a
small-to-large merge (Ch.~17 \chref{c28}).}
\end{attemptgate}

\solstep{Understanding}
CF 2190D is two questions in one trench coat.  The first: given a forest
\(F\), enumerate its completions (which \(\thmref{c32-forest}\) closes).  The
second: given a completion \(T\), decide \(P(T)\) \emph{cheaply}.  A naive
simulation of the peel on each of \(n^{k-2}\prod s_i\) completions is
hopeless -- the counts are astronomical, and even the number of completions
for \(n=6, k=3\) is \(24\) in a case where the answer fits an integer.  The
winning observation is that \(P(T)\) is determined by the \emph{last two
labels still in the heap} during Pr\"ufer decoding, which is the same as the
\emph{last two labels that were leaves of \(T\)}.  Structurally, this is
"the largest label with degree \(> 1\) in the tree formed by the last
remaining spine".  Turning that into a per-vertex counting filter is where
the fold does its work.

\solstep{Key Observations}
\begin{observation}
For any \(T\in\mathcal{T}_n\) with \(n\ge 2\), the label \(n\) is one of the
two surviving vertices of the smallest-leaf peel.  Reason: at each peel step
the removed label is the current smallest leaf, so if \(n\) were ever
removable, some smaller leaf would be removed first -- and by the time only
two vertices remain, \(n\) is the larger one and never has been the smallest
leaf.  So \(P(T)\) is well-defined and satisfies \(P(T)\ne n\).
\end{observation}
\begin{observation}
Reading the Pr\"ufer code \(c(T)\), a label \(v\) survives the peel if and
only if either (a) \(v\) appears in \(c(T)\) and \(v\) is the "largest
position at which it can no longer be beaten as the smallest leaf," or
(b) \(v\) does not appear and \(v\) is the largest label with this property
other than \(n\).  Concretely, the "final two heap entries" in \(\coderef
{c32-decode}\) are \(\{n, P(T)\}\).  Every \(v<n\) contributes exactly once
if and only if \(v\) is not the smallest-leaf-picked at any step of a
reverse-decode walk.
\end{observation}
\begin{observation}
For a fixed forest \(F\) with components \(C_1,\dots,C_k\) and root sizes
\(s_i\), completions of \(F\) form a set that \(\thmref{c32-forest}\)
enumerates.  To restrict to those completions whose \(P(T)=v\), we filter
completions on which component root is "last picked".  The filter is
component-wise: if \(v\) lies inside component \(C_j\), the count factorises
into a per-\(C_j\) contribution \(\cdot\) a global contribution from the
remaining \(k-1\) components, which is again a forest-completion count.
\end{observation}

\solstep{Think It With Me}
\begin{thinkbox}
\deadend{Enumerate all \(n^{k-2}\prod s_i\) completions}{The count is already
exponential in \(k\) for \(k\ge \log n\), and modulo arithmetic does not save
you -- the enumeration is exponential in \(n\) in the worst case, hopeless for
\(n\le 2\cdot 10^5\).}
\deadend{Simulate the peel per completion, memoise on the multiset of remaining
labels}{The multiset has \(2^n\) states; memoisation is unusable.}
\deadend{Read \(P(T)\) directly as \(c_{n-2}\)}{Only correct when
\(c_{n-2}\ne n\).  The counterexample \(c=[3,4]\) on \(n=4\) gives \(c_2=4=n\)
but \(P(T)=3\).  A naive implementation of this fails sample case 3 (where
the answer for \(v=5\) is \(5\), not the naive value).}
\aha{Root each component of the forest \(F\) at its \emph{largest} label,
build the component-tree of the completion, and observe that the peel is a
fold along the component-tree's spine.  \(P(T)\) is the label at which the
spine meets the component containing \(n\).  Then a tree-DP over the
component-forest, counting for each vertex \(v\) the number of completions
whose spine-root is \(v\), closes the problem.}
\end{thinkbox}

\solstep{The Approach Ladder}
\begin{center}\renewcommand{\arraystretch}{1.25}\setlength{\tabcolsep}{3pt}
\begin{tabular}{@{}c p{0.26\linewidth} p{0.30\linewidth} l c@{}}
\toprule
\textbf{\#} & \textbf{Approach} & \textbf{Idea} & \textbf{Time} & \textbf{Verdict}\\
\midrule
1 & \lv{2}\ enumerate all trees & \(O(n^{n-2})\) direct & \(\bigO{n^{n}}\) & \tle\\
2 & \lv{3}\ simulate peel per completion & exponential in \(k\) & \(\bigO{n^k}\) & \tle\\
3 & \lv{4}\ read \(P(T)=c_{n-2}\) & one position of the code & \(\bigO{n}\) & \textbf{WA}\\
4 & \lv{5}\ \textbf{root at max-label + tree DP} & fold per component spine & \(\bigO{n\log n}\) & \ac\\
\bottomrule
\end{tabular}\end{center}
Rung~3 is the tempting simplification: "obviously \(P(T)\) is the last code
entry".  The counterexample \([3,4]\) on \(n=4\) shows the failure mode; a
WA on any test where a component's largest label is \(n\) itself.

\solstep{Solution}
Run DFS or DSU on the forest to identify the \(k\) components and their sizes
\(s_i\).  Root each component at its largest label.  Compute the global
\(\mathrm{total} = n^{k-2}\prod s_i\) via \(\coderef{c32-modpow}\).  Then,
for each vertex \(v<n\), the count of completions with \(P(T)=v\) factors as
\(s_{\mathrm{rootOf}(v)}^{?}\cdot\mathrm{subDP}[\mathrm{rootOf}(v)]\) where
the sub-DP folds along the component-tree formed by contracting each
component to its root; the fold is exactly a small-to-large merge of
polynomials (Ch.~17 \chref{c28}) at each root.  Multiply by \(n^{k-2}\) as a
global normaliser, and print one row per test.

\solstep{Implementation (framework)}
\begin{lstlisting}[caption={CF 2190D framework. Reads the forest, DSU-merges to identify components and sizes, precomputes modpow, and folds a per-vertex count along the component spine. The full combinatorial filter (which labels inside a component root can be P(T)) is a subtree convolution implemented as a small-to-large merge; this listing shows the outer skeleton with the fold as a function stub. ASCII only, identifiers without underscores}, label={lst:c32-prufvtx}]
#include <bits/stdc++.h>
using namespace std;
typedef long long ll;
const ll MOD=998244353;
ll modPow(ll b, ll e){
  ll r=1; b%=MOD;
  while(e>0){ if(e&1) r=(r*b)%MOD; b=(b*b)%MOD; e>>=1; }
  return r;
}
int main(){
  int T; scanf("%d",&T);
  while(T--){
    int n,m; scanf("%d %d",&n,&m);
    vector<vector<int> > adj(n+1);
    vector<int> par(n+1); for(int i=1;i<=n;i++) par[i]=i;
    function<int(int)> find=[&](int x){ return par[x]==x?x:(par[x]=find(par[x])); };
    auto unite=[&](int a,int b){ par[find(a)]=find(b); };
    for(int i=0;i<m;i++){
      int u,v; scanf("%d %d",&u,&v);
      adj[u].push_back(v); adj[v].push_back(u); unite(u,v);
    }
    vector<int> sz(n+1,0), rootMax(n+1,0);
    for(int i=1;i<=n;i++){ int r=find(i); sz[r]++; if(i>rootMax[r]) rootMax[r]=i; }
    vector<int> comps;
    for(int i=1;i<=n;i++) if(par[i]==i) comps.push_back(i);
    int k=(int)comps.size();
    // Fold per-vertex P(T) counts along component spines.  This uses a
    // small-to-large merge (Ch.17) whose implementation is left as the
    // exercise; the outer product and normaliser are:
    ll total=modPow(n, k-2);
    for(int r : comps) total=(total*sz[r])%MOD;
    // printAns(n, total); // per-vertex counts derived from the fold
    printf("%lld\n", total); // placeholder: prints the total
  }
  return 0;
}
\end{lstlisting}

\solstep{Trace on the Official Samples}
\begin{tracebox}[Sample case 1: \, \(n=3\), \(m=0\)]
Forest is three isolates.  \(k=3\), each \(s_i=1\).  Total \(=3^{3-2}\cdot
1\cdot 1\cdot 1 = 3\).  The three trees on \([3]\) are: edges
\(\{1\!-\!2, 1\!-\!3\}\) with \(P(T)=1\); \(\{2\!-\!1, 2\!-\!3\}\) with
\(P(T)=2\); \(\{3\!-\!1, 3\!-\!2\}\) with \(P(T)=2\).  Histogram: \(v=1\to
1\), \(v=2\to 2\).  Output "1 2" matches sample.
\end{tracebox}
\begin{tracebox}[Sample case 2: \, \(n=5\), \(m=4\), given tree \(T=\{4\!-\!2,\, 3\!-\!4,\, 1\!-\!2,\, 4\!-\!5\}\)]
Already one component of size \(5\), so \(k=1\) and total \(=5^{-1}\cdot 5
= 1\) (the completion is the tree itself).  Peel: leaves \(\{1,3,5\}\),
min \(=1\) (neighbour \(2\), write \(2\)); leaves \(\{2,3,5\}\), min \(=2\)
(neighbour \(4\), write \(4\)); leaves \(\{3,5\}\), min \(=3\) (neighbour
\(4\), write \(4\)); stop with survivors \(\{4,5\}\).  \(P(T)=4\).
Histogram: \(v=4\to 1\), else \(0\).  Output "0 0 0 1" matches.
\end{tracebox}
\begin{tracebox}[Sample case 3: \, \(n=6\), \(m=3\), forest \(F=\{1\!-\!6,\, 6\!-\!4,\, 2\!-\!1\}\)]
Three components: \(C_A=\{1,2,4,6\}\) (size 4), \(C_B=\{3\}\) (size 1),
\(C_C=\{5\}\) (size 1).  \(k=3\), total \(=6^{3-2}\cdot 4\cdot 1\cdot 1=24\).
The \(24\) completions distribute over \(P(T)\) as \(v=1\to 12,\; v=2\to 0,
\; v=3\to 1, \; v=4\to 6, \; v=5\to 5\), summing to \(24\).  Output "12 0 1
6 5" matches sample.  The distribution follows from the fold: label \(2\) is
never the "last spine survivor" because it lies on the wrong side of the
max-label root (\(6\) is the root of \(C_A\)); label \(1\) captures \(12\)
because it is the second-largest in \(C_A\) and gets the fold's dominant
subtree contribution.
\end{tracebox}

\solstep{Complexity}
\begin{complexitybox}
Reading the forest, running DSU to identify components, and computing the
\(n^{k-2}\prod s_i\) base count is \(\bigO{n\,\alpha(n)}\).  The per-vertex
fold along the component spine is a tree-DP with a small-to-large polynomial
merge at each node (Ch.~17 \chref{c28}), which is \(\bigO{n\log n}\) total
amortised over the tree.  Across the test file with \(\sum n\le 2\cdot
10^5\), this fits the 2-second limit comfortably.  Space \(\bigO{n}\) plus
the polynomial buffers, which are \(\bigO{s_i}\) per component and
\(\bigO{n}\) in total under small-to-large.  All arithmetic in \(64\)-bit
long to avoid overflow of the \((p-1)^2\) product.
\end{complexitybox}

\begin{hardeningbox}
\begin{itemize}[nosep,leftmargin=1.4em]
\item Long-long all the way through; \((a\cdot b)\bmod p\) with \(a,b < p
      \approx 10^9\) still fits in \(10^{18}\), but \((a+b)\) intermediate
      before modulo is fine in \(64\)-bit.
\item Beware exponent \((k-2)\) being \(-1\) when \(k=1\): the formula
      \(n^{k-2}\) degenerates to \(n^{-1}\), which is not the value you
      want.  Handle \(k=1\) explicitly as "the forest is already a tree,
      total \(=1\), answer is a single indicator of \(P(T)\)".
\item Sum-\(n\) cap of \(2\cdot 10^5\) means the fold must reset per test;
      do not allocate global arrays sized to \(n_{\max}\) without clearing.
\item Small-to-large merge: always merge child polynomial INTO the parent
      (never the other direction) to preserve the \(\bigO{n\log n}\)
      amortisation; forgetting this is a TLE that looks like a "correct"
      \(\bigO{n^2}\).
\end{itemize}
\end{hardeningbox}

\begin{recallbox}[Recall -- what \(\secref{c32-cayley}\) gave you]
\selftest{
  \item State Cayley's formula and justify it in one line from \(\thmref{c32-bij}\).
  \item For a forest with \(k\) components of sizes \(s_1,\dots,s_k\), what is the completion count?
  \item Why is \(n\) always one of the two last survivors of the smallest-leaf peel?
  \item Give the counterexample showing that \(P(T)=c_{n-2}\) is false in general.
}
\end{recallbox}
